\section{Introduction}
\label{sec:intro}

Since the work of Pareto, efficiency has been the cornerstone of multicriteria analysis: an alternative is Pareto efficient if no criterion can be improved without worsening another. This is a safety property that removes dominated alternatives, but it is decisionally incomplete. For $m$ conflicting criteria the efficient set is typically large, leaving the decision maker (DM) with no principled selection rule. The fundamental issue remains that the Pareto front is a safety filter, not a complete decision rule.

LexUR turns the output from a set into a certified recommendation. The DM declares a monotone probe family and normalization bounds. Each alternative receives normalized regrets over that family. These regrets are lexicographically sorted, and selection follows a lexicographic minimax rule. This yields a certified stability class accompanied by an explanation of which questions remain most severely binding. Figure~\ref{fig:paretocert} contrasts the safety filter with the regret certificate.

% From a Pareto front (safety filter) to a regret certificate (decision)
\begin{figure}[t]
\centering
\begin{tikzpicture}
% ---------- left panel: efficient front ----------
\begin{scope}
\draw[->, line width=0.6pt] (0,0) -- (4.2,0) node[right, font=\scriptsize] {$f_1$};
\draw[->, line width=0.6pt] (0,0) -- (0,4.0) node[above, font=\scriptsize] {$f_2$};
% convex-ish front of points
\foreach \x/\y in {0.4/3.5, 0.7/2.7, 1.1/2.1, 1.6/1.6, 2.2/1.2, 3.0/0.85, 3.6/0.6}{
  \fill[lurslate] (\x,\y) circle (2.0pt);
}
\node[lurslate, font=\scriptsize, align=center] at (2.4,3.2) {efficient set:\\\emph{which one?}};
\node[font=\scriptsize\bfseries, lurslate] at (2.1,-0.55) {safety filter};
% highlight chosen point
\fill[lurcrimson] (1.6,1.6) circle (3.0pt);
\draw[lurcrimson, line width=0.7pt] (1.6,1.6) circle (5.5pt);
\end{scope}

% ---------- arrow ----------
\node[flow, single arrow, draw=lurcrimson, fill=lurcrimson!10, minimum height=12mm,
      single arrow head extend=1.4mm, inner sep=3pt, font=\scriptsize, align=center]
      at (5.2,2.0) {LUR};

% ---------- right panel: certificate ----------
\begin{scope}[xshift=6.6cm, yshift=0.2cm]
\draw[->, line width=0.6pt] (0,0) -- (0,3.6) node[above, font=\scriptsize] {};
\node[font=\scriptsize, rotate=90] at (-0.45,1.7) {sorted disappointment};
\foreach \i/\h/\lab in {0/2.9/$q_{[1]}$, 1/2.0/$q_{[2]}$, 2/1.3/$q_{[3]}$, 3/0.8/$q_{[4]}$, 4/0.4/$q_{[5]}$}{
  \pgfmathsetmacro\xx{0.35+\i*0.65}
  \ifnum\i<2
    \fill[lurcrimson] (\xx,0) rectangle ++(0.45,\h);
  \else
    \fill[lurslate!55] (\xx,0) rectangle ++(0.45,\h);
  \fi
  \node[font=\scriptsize] at (\xx+0.22,-0.28) {\lab};
}
\draw[line width=0.6pt] (0,0) -- (3.7,0);
\node[lurcrimson, font=\scriptsize\bfseries, align=center] at (1.9,3.3) {binding probes};
\node[font=\scriptsize\bfseries, lurslate] at (1.9,-0.75) {decision certificate};
\end{scope}
\end{tikzpicture}
\caption{From a front to a certificate. The efficient set (left) removes dominated
options but leaves the choice open. LUR returns one alternative together with its
sorted regret certificate (right), whose largest entries---the \emph{binding
probes}---explain where the recommendation remains most fragile.}
\label{fig:paretocert}
\end{figure}


Our contributions are unified across four dimensions. Theoretically, we define the LexUR order and prove existence, completeness, Pareto compatibility, and stability under bounded perturbations. Algorithmically, we provide an exact finite-set procedure using sequential sorts. Empirically, we validate LexUR on a broadened benchmark suite, demonstrating competitive tail performance against strong robust baselines. Interpretably, we introduce an adaptive correlation-clustering probe family that handles redundancy and produces an auditable regret certificate for the final choice.

Robustness in this framework is strictly relative to the declared probe family, not all possible preferences. Furthermore, point normalization may be sensitive when ideal and nadir bounds are uncertain. LexUR is not parameter-free, but its assumptions are declared, auditable, and separated from the final recommendation.

Section~\ref{sec:related} positions the method. Sections~\ref{sec:core} to~\ref{sec:probes} develop the core theory, computation, and probe design. Sections~\ref{sec:ext} and~\ref{sec:exp} test extensions and empirical performance, and Section~\ref{sec:disc} closes with implications.
